\section{Alignment e ingeniería de datos: adaptar la conducta no es generar conocimiento de la nada}

Una vez entrenada la base del lenguaje, SFT y preference optimization determinan cómo responde el modelo. El ponente agrupa esta parte con data engineering porque la variable decisiva de alignment no es el nombre del algoritmo, sino la calidad, la procedencia y la cobertura de las instruction, las answer y los preference pair.

\subsection{SFT: las respuestas de alta calidad requieren domain experts}

SFT (Supervised Fine-Tuning) sigue siendo la cross entropy teacher-forced habitual. Dado un par instruction--response $(x,y)$, el objetivo es
\[
\mathcal{L}_{\text{SFT}}=-\sum_{t=1}^{|y|}\log \pi_\theta(y_t\mid x,y_{<t}).
\]
La dificultad está en los datos: si se quiere que el modelo escriba explicaciones de código confiables, quien anota debe ser capaz de escribir explicaciones de código confiables, y no limitarse a resolver preguntas de opción múltiple de bajo costo.

\lecturefigure{slide-10-sft.jpg}{SFT: expert annotation, distillation y el problema weak-to-strong}{Grabación oficial de Stanford Online 00:27:01}

\begin{knowledgebox}{Los dos límites de los teacher-generated data}
Un modelo más fuerte puede generar instruction data con rapidez, pero el estudiante hereda los sesgos, el estilo y los puntos ciegos del teacher; si el objetivo es competir de forma independiente, depender de un teacher cerrado trae además problemas de licencias, reproducibilidad y data provenance. Weak-to-Strong estudia si, bajo supervisión débil, se puede superar al supervisor; no equivale a «copiar las teacher outputs basta para superar necesariamente al teacher».
\end{knowledgebox}

\teachervoice{00:25:00--00:28:24, el ponente insiste en que SFT no es crowdsourcing ordinario. Las respuestas de alta calidad requieren personas que de verdad conozcan el domain, y esto importa más que reducir la cantidad de datos a una cifra.}

\subsection{RLHF y DPO: el algoritmo se simplifica, pero el supuesto de preferencias no desaparece}

SFT enseña al modelo a imitar la respuesta objetivo, pero no compara de forma explícita dos salidas que «parecen respuestas» por igual. Preference optimization responde entonces la siguiente pregunta: cuando chosen y rejected solo difieren en utilidad, estilo o seguridad, cómo convertir un juicio por pares en una policy update. La diferencia entre RLHF basado en PPO y DPO está en la ruta de optimización, no en si hace falta definir las preferencias.

RLHF al estilo PPO suele entrenar primero un reward model $r_\phi(x,y)$ y después optimizar la policy con KL regularization:
\[
\max_\theta\ \mathbb{E}_{y\sim\pi_\theta(\cdot\mid x)}[r_\phi(x,y)]
-\beta D_{\mathrm{KL}}\!\left(\pi_\theta\,\|\,\pi_{\text{ref}}\right).
\]
DPO optimiza directamente con el pair chosen/rejected $(y^+,y^-)$
\[
\mathcal{L}_{\text{DPO}}
=-\log\sigma\!\left(\beta\left[
\log\frac{\pi_\theta(y^+\mid x)}{\pi_{\text{ref}}(y^+\mid x)}
-\log\frac{\pi_\theta(y^-\mid x)}{\pi_{\text{ref}}(y^-\mid x)}
\right]\right).
\]

\lecturefigure{slide-11-rlhf.jpg}{RLHF y DPO: reward model, la dificultad de PPO y el pairwise objective}{Grabación oficial de Stanford Online 00:28:24}

\begin{warningbox}{DPO no elimina el reward modeling}
DPO no entrena de forma explícita un reward model independiente, pero los preference pairs siguen definiendo de manera implícita «qué es mejor». Que el pair sea on-policy o no, que los annotator sean consistentes, que los prompt cubran el uso real y que chosen/rejected difieran solo en estilo: todo ello modifica la reward geometry que se aprende.
\end{warningbox}

\teachervoice{00:28:24--00:30:00, el ponente considera que PPO puede ser muy potente cuando el reward model es lo bastante bueno, pero su implementación es difícil; presenta DPO como una alternativa más fácil de usar y, al mismo tiempo, advierte de palabra sobre la pair quality y el carácter on-policy.}

\subsection{Data, algorithm y architecture pueden codificar un bias similar}

Esta slide contiene la metodología de investigación más importante de toda la clase. El multi-hop question answering de CogQA llegó a usar una graph neural network, y otra línea usó MCTS; cuando un GPT de long-context y chain-of-thought pueden colocar de una sola vez en el context todos los documents relevantes, la misma tarea puede reformularse como un problema a nivel de data/context.

\lecturefigure{slide-12-data-algorithm-architecture.jpg}{Intercambiabilidad entre datos, algoritmo y arquitectura: de CogQA al razonamiento con long-context}{Grabación oficial de Stanford Online 00:34:16}

\begin{importantbox}{Tres niveles para implementar el mismo sesgo inductivo}
\begin{enumerate}
  \item \textbf{Architecture:} incorporar graph structure, memory o routing en la estructura de la red.
  \item \textbf{Algorithm:} organizar el cómputo en inference mediante search, MCTS, retrieval o un verifier.
  \item \textbf{Data/context:} colocar la evidencia intermedia, las rutas fallidas o los documents completos en los ejemplos de entrenamiento y en el contexto, para que un modelo general aprenda la procedure.
\end{enumerate}
Cuanto más cerca del nivel de los datos, más general suele ser la solución, pero exige un modelo más fuerte, un context más largo y ejemplos de mayor calidad.
\end{importantbox}

\begin{lstlisting}[caption={Marco de decisión para elegir entre data, algorithm o architecture}]
if behavior_is_narrow_and_examples_are_available:
    encode_with_data_and_evaluation()
elif inference_needs_verifiable_search:
    add_algorithmic_search_or_tools()
elif bottleneck_is_general_and_repeats_across_tasks:
    consider_architecture_change()
always_measure(compute, latency, coverage, failure_modes)
\end{lstlisting}

\teachervoice{00:30:00--00:34:20, el ponente defiende el data work: cleaning, filtering y synthesizing no son «trabajo de bajo nivel», sino la ingeniería de investigación más central de las empresas de modelos actuales.}

\teachervoice{01:10:30--01:12:30, en la sesión de Q\&A añadió un matiz: muchas conductas concretas pueden resolverse con data, pero si alguien encuentra una architecture update que mejore de forma general la capacidad de ajuste, sigue teniendo mucho valor.}

\subsection{Resumen de la sección}

SFT, RLHF y DPO no pueden discutirse al margen de la calidad de los datos. En términos más generales, un mismo inductive bias puede incorporarse en architecture, algorithm o data/context; el nivel más sencillo no es necesariamente el más barato, porque puede exigir más compute, un context más largo y una evaluation más confiable.
